% Part: proof-theory
% Chapter: sequent-calculus
% Section: rules-proofs

\documentclass[../../../include/open-logic-section]{subfiles}

\begin{document}

\olfileid{pt}{seq}{rp}

\olsection{في القواعد و\usetoken{P}{proof}}

نقدم أولًا ما نحتاج إليه من التعريفات، ونضبط ألفاظها.

فمن ذلك قاعدتا~\Log{G3c} اللتان تتيحان~!!{proving} متتاليات تشتمل على
$!A\land !B$ من متتاليات تشتمل على~$!A$ و$!B$:
\[
\Axiom$ !A, !B, \Gamma \fCenter \Delta$
\RightLabel{\LeftR{\land}}
\UnaryInf$ !A \land !B, \Gamma \fCenter \Delta$
\DisplayProof
\qquad
\Axiom$\Gamma \fCenter \Delta, !A$
\Axiom$ \Gamma \fCenter \Delta, !B$
\RightLabel{\RightR{\land}}
\BinaryInf$ \Gamma \fCenter \Delta, !A \land !B $
\DisplayProof
\]
فالواقع فوق الخط من المتتاليات يسمى \emph{مقدمات}، والواقع تحته يسمى
\emph{النتيجة}. وأما~!!{formula}s التي في~$\Gamma$ و$\Delta$ في كل قاعدة
فهي \emph{السياق}، وتسمى كذلك~!!{formula}s \emph{الجانبية}. وما كان في
النتيجة من~!!{formula} خارج السياق فهو~!!{formula} \emph{الرئيسة}؛ وما
كان في المقدمات من~!!{formula}s خارج السياق فهو~!!{formula}s
\emph{الفعالة}، ويقال لها أيضًا \emph{المساعدة}.

وأما قواعد الكميات ففيها مزيد دقة. فقواعد~$\lforall$ في~\Log{G1c} مثلًا:
\[\Axiom$ !A(t), \Gamma \fCenter \Delta$
\RightLabel{\LeftR{\lforall}}
\UnaryInf$ \lforall[x][!A(x)],\Gamma \fCenter \Delta$
\DisplayProof
\qquad
\Axiom$ \Gamma \fCenter \Delta, !A(a) $
\RightLabel{\RightR{\lforall}}
\UnaryInf$ \Gamma \fCenter \Delta, \lforall[x][!A(x)]$
\DisplayProof
\]
في قاعدة~\LeftR{\lforall} تكون~!!{formula} الفعالة~$!A(t)$، ويكون
$t$ حدًا مغلقًا، أي خاليًا من المتغيرات. ولا يصح هذا الاستدلال---أي لا
يطابق مخطط~\LeftR{\lforall}---إلا إذا وجدت~!!{formula}~$!A$ ومتغير~$x$
وحد مغلق~$t$، وكانت~!!{formula} الرئيسة في النتيجة
$\lforall[x][!A]$، والـ!!{formula} الفعالة في المقدمة
$\Subst{!A}{t}{x}$. وعلى هذا النحو تعمل~\RightR{\lforall}، إلا أن
!!{formula} المقدمة الفعالة يجب أن تكون~$\Subst{!A}{a}{x}$، حيث~$a$
!!a{constant}. وهذا الثابت هو \emph{المتغير المميَّز} لاستدلال
\RightR{\lforall}. وشرط صحة الاستدلال ألا تحتوي المتتالية السفلى~$a$؛
أي ألا يقع~$a$ في~$\Gamma$ أو~$\Delta$ أو~$!A$. وهذا بعينه
\emph{شرط المتغير المميَّز}.

يتألف نظام المتتاليات، كـ\Log{G1c} و\Log{G3c}، من قواعد معينة بالمخططات
المتقدمة، ومن متتاليات معينة مخططيًا تتخذ بديهيات. وأما~!!^{proof}s فيه
فأشجار منتهية من المتتاليات، تنشأ استقرائيًا من البديهيات بقواعد
الاستدلال: كل ورقة فيها بديهية، وكل متتالية سوى الأوراق تترتب، بقاعدة من
قواعد النظام، على المتتاليات التي تعلوها مباشرة.

\begin{defn}\ollabel{defn:proof}
تُنشأ~\emph{!!^{proof}s} نظام المتتاليات استقرائيًا، وهي أشجار منتهية من
المتتاليات، على الوجه الآتي:
\begin{enumerate}
  \item\ollabel{defn:prf:axiom} البديهية الواحدة~!!a{proof}.
  \item\ollabel{defn:prf:one-prem} إذا كان~$\pi_1$~!!a{proof} خاتمته
  المتتالية~$S_1$، وكانت~$S$ تترتب على~$S_1$ بقاعدة ذات مقدمة واحدة~$R$، فإن
  \[
  \AxiomC{}
  \RightLabel{$\pi_1$}
  \DeduceC{$S_1$}
  \RightLabel{$R$}
  \UnaryInfC{$S$}
  \DisplayProof
  \]
  يكون~!!a{proof}.
  \item\ollabel{defn:prf:two-prem} وإذا كان~$\pi_1$ و$\pi_2$~!!{proof}s
  خاتمتاهما~$S_1$ و$S_2$، وكانت~$S$ مترتبة على~$S_1$
  و~$S_2$ بقاعدة ذات مقدمتين~$R$، فإن
  \[
  \AxiomC{}
  \RightLabel{$\pi_1$}
  \DeduceC{$S_1$}
  \AxiomC{}
  \RightLabel{$\pi_2$}
  \DeduceC{$S_2$}
  \RightLabel{$R$}
  \BinaryInfC{$S$}
  \DisplayProof
  \]
  يكون~!!a{proof}.
\end{enumerate}
\end{defn}

يجعل~\Olref{defn:proof}~!!{proof}s أشجارًا من المتتاليات، ونتصور نموها
إلى أعلى. فعقدها العليا، وهي الأوراق، بديهيات تسمى \emph{متتاليات
ابتدائية}؛ والمتتالية السفلى هي \emph{المتتالية النهائية}. فإذا كان
$\pi$~!!a{proof} خاتمته~$S$، كتبنا كذلك~$\pi\Proves S$. وإذا وجد
!!a{proof} للمتتالية~$S$، كتبنا~$\Proves S$، وربما أثبتنا اسم النظام
الذي فيه البرهان إلى يسار~$\Proves$؛ ومثاله
$\Log{G1c}\Proves !A,!B\Sequent !A\land !B$. وكل متتالية في~!!a{proof}
سوى الابتدائية \emph{نتيجة} لاستدلال بقاعدة~$R$. والمتتالية أو
المتتاليتان اللتان تعلوان متتالية ما مباشرة في~!!a{proof}، أي أبوا
العقدة، هما \emph{مقدمات} ذلك الاستدلال.

والـ!!{proof}s الداخلة في البناء الاستقرائي لـ!!a{proof} تسمى
\emph{!!{proof}s جزئية} منه؛ وتسمى~!!{proof}s الجزئية المنتهية بمقدمات استدلال بعينه
\emph{!!{proof}s جزئية مباشرة} لذلك الاستدلال.

وكثيرًا ما نقيس تعقيد~!!{proof}s بعدد طبيعي. وقد نعد المتتاليات في
!!a{proof}، غير أن الأنفع في الغالب هو~!!{height}~!!a{proof}: أي أكبر
عدد لخطوات الاستدلال، أو للحواف، في فرع يمتد من المتتالية النهائية إلى
متتالية ابتدائية. وله تعريف استقرائي أيضًا.

\begin{defn}
يُعرَّف~\emph{!!{height}}~$\pheight{\pi}$ لـ!!a{proof}~$\pi$ استقرائيًا
على الوجه الآتي:
\begin{enumerate}
  \item إن اقتصر~$\pi$ على بديهية واحدة، فـ$\pheight{\pi}=0$.
  \item وإن ختم~$\pi$ استدلال ذو مقدمة واحدة، وكان~$\pi_1$ هو
  الـ!!{proof} الجزئي المباشر، فـ$\pheight{\pi}=\pheight{\pi_1}+1$.
  \item وإن ختم~$\pi$ استدلال ذو مقدمتين، وكان~$\pi_1$ و$\pi_2$
  هما الـ!!{proof}s الجزئية المباشرة، فـ
  $\pheight{\pi}=\max(\pheight{\pi_1},\pheight{\pi_2})+1$.
\end{enumerate}
ومتى كان للمتتالية~$S$~!!a{proof} ارتفاعه~$\le n$، كتبنا~$\Proves[n]S$.
\end{defn}

\begin{ex}
كل متتالية على الصورة~$!C,\Gamma\Sequent\Delta,!C$ بديهية في~\Log{G3c}،
إذا كانت~$!C$ ذرية. ولتكن~$!D$ و$!E$ جملتين ذريتين؛ فتكون
$!D,!E\Sequent !D$ و$!D,!E\Sequent !E$ صورتين خاصتين من البديهية
$!C,\Gamma\Sequent\Delta,!C$. فإذا وضعنا، مثلًا،~$!C\ident !E$ و
$\Gamma=\{!D\}$ و$\Delta=\emptyset$، نتجت بعينها المتتالية
$!D,!E\Sequent !E$. ولأن السياقات متعددات مجموعات، فإن
$!D,!E\Sequent !E$ و$!E,!D\Sequent !E$ متتالية واحدة.

ولما كانت~$!D,!E\Sequent !D$ صورة من بديهيات~\Log{G3c}، كانت وحدها
!!a{proof} في~\Log{G3c}؛ وكذلك~$!D,!E\Sequent !E$ بمقتضى
\olref{defn:proof}\olref{defn:prf:axiom}. ولنسمهما~$\pi_1$ و$\pi_2$.
ثم يجيز~\olref{defn:prf:two-prem} أن نركب من هذين~!!{proof}s واحدًا
جديدًا~($\pi_3$) بقاعدة~\RightR{\land} ذات المقدمتين:
\[
\Axiom$!D, !E \fCenter !D$
\Axiom$!D, !E \fCenter !E$
\RightLabel{\RightR{\land}}
\BinaryInf$!D, !E  \fCenter !D \land !E $
\DisplayProof
\]
ومن~$\pi_3$ ننتقل إلى~!!{proof}~$\pi_4$:
\[
\Axiom$!D, !E \fCenter !D$
\Axiom$!D, !E \fCenter !E$
\RightLabel{\RightR{\land}}
\insertBetweenHyps{\hspace{5ex}}
\BinaryInf$!D, !E  \fCenter !D \land !E $
\LeftSubproofLabel{$\pi_3$}
\RightLabel{\LeftR{\land}}
\UnaryInf$!D \land !E  \fCenter !D \land !E$
\RightSubproofLabel{$\pi_4$}
\DisplayProof
\]
وذلك باستعمال قاعدة~\LeftR{\land} ذات المقدمة الواحدة، وفق
\olref{defn:proof}\olref{defn:prf:one-prem}. فالـ!!{proof} الجزئي
المباشر لاستدلال~$\pi_4$ الأخير هو~$\pi_3$؛ والـ!!{proof}s الجزئية
المباشرة لاستدلال~\RightR{\land} هما~$\pi_1$ و$\pi_2$. وأما
الـ!!{proof}s الجزئية لـ$\pi_4$ فتشمل~$\pi_1$ و$\pi_2$ و$\pi_3$
و$\pi_4$ نفسه. ثم إن~$\pheight{\pi_1}=\pheight{\pi_2}=0$ و
$\pheight{\pi_3}=1$ و$\pheight{\pi_4}=2$. فثبت أن
$\Log{G3c}\Proves[2] !D\land !E\Sequent !D\land !E$ لكل~$!D$ و$!E$
ذريتين.
\end{ex}

\end{document}
